\documentclass[10pt,a4paper]{amsart}
\usepackage[margin=19mm]{geometry}
\usepackage{amsmath,amssymb,mathtools}
\usepackage[scr=rsfs,cal=euler]{mathalfa}
\usepackage{xfrac}
\usepackage[unicode,hidelinks,bookmarks=false]{hyperref}
\newcommand{\ie}{i.e.\ }
\newcommand{\cf}{cf.\ }
\newcommand{\resp}{respectively\ }
\newcommand{\Subsection}{\subsection}
\providecommand{\nobreakdash}{\nobreak\hbox{-}}
\setlength{\emergencystretch}{2em}
\multlinegap0pt
\usepackage{mathtools}
\usepackage{amssymb}
\usepackage{bm}
\usepackage{xcolor}
\usepackage[normalem]{ulem}
\usepackage{cancel}
\newtheorem{problem}{Problem}
\newtheorem{proposition}{Proposition}
\newcommand{\IR}{\mathbb{R}}
\newcommand{\IV}{\mathcal{V}}
\newcommand{\MS}[1]{{\color{red} #1}}
\DeclareMathOperator*{\argmin}{arg\,min}
\newcommand{\cH}{\mathcal{H}}
\newcommand{\cV}{\mathcal{V}}
\newcommand{\cVrb}{\cV^\rb_R}
\newcommand{\dotp}[2]{\left ( #1,#2 \right )}
\newcommand{\dual}[2]{\left \langle #1,#2 \right \rangle}
\newcommand{\norm}[1]{\left \lVert #1 \right \rVert}
\newcommand{\rb}{{(\mathrm{rb})}}
\newcommand{\sfA}{\mathsf{A}}
\newcommand{\vrb}{v^\rb_R}
\title{A Model Order Reduction Method for Seismic Applications Using the Laplace Transform}
\author{Fernando Henriquez, Matthias Schlottbom}
\date{Source: arXiv:2602.15517v1 (CC BY 4.0)}
\begin{document}
\maketitle

\noindent\emph{Source and licence.} A Model Order Reduction Method for Seismic Applications Using the Laplace Transform. Version arXiv:2602.15517v1 (CC BY 4.0), distributed by arXiv under the Creative Commons Attribution 4.0 International licence, http://creativecommons.org/licenses/by/4.0/, as recorded in the arXiv licence metadata for this exact version and shown on the abstract page https://arxiv.org/abs/2602.15517v1. Full licence notice: \texttt{licenses/arxiv-2602.15517-CC-BY-4.0.txt}.\par\smallskip

\noindent\textit{Editorial scope.} Counted unit: the article's proposition eq:best\_approximation (source0.tex lines 4313--4354). The article's complete original proof of it is delivered with it (source0.tex lines 4356--4698, one \texttt{proof} environment of 343 lines). No mathematics is written, repaired or re-derived: the statement and the proof are the article's own lines, in the article's own notation, and no retained line is substituted or re-flowed.  Retained uncounted as the source-local context the proof needs: lines 490--494; lines 962--988; lines 964--985; lines 4261--4277; lines 4289--4303; lines 4305--4310.  Nothing was omitted from any retained span, so the package carries no omission record.  No retained line contains a citation, so the wrapper carries no bibliography and no bibliography entry.  No retained line contains a figure environment, an image-inclusion command or drawing code, so no image is delivered and the package declares no asset dependency.  Editorial formatting only, in addition: an AMS \texttt{amsart} preamble that restates the article's own declarations and macros used by the retained lines, namely the environments \texttt{problem}, \texttt{proposition} on separate counters, together with the standard packages those lines need.  The retained environments print in this document as Problem 1, Proposition 1 in order of appearance, whereas the article prints the same items with the numbering of its own full document.  The complete native source of the article as distributed in the arXiv e-print for this exact version is bundled unchanged as \texttt{sources/194804/source0.tex}.
\begin{align}
    \partial_t^2 u(t) + \mathsf{A} u(t) &= f(t), \quad t>0, \label{eq:wave}\\
    u(0) &= 0, \label{eq:ic1}\\
    \partial_t u(0) &= 0. \label{eq:ic2}
\end{align}

\begin{problem}[Reduced semi-discrete problem] \label{pr:sdpr}
We seek $u^{\normalfont\textrm{(rb)}}_{R,M} \in \mathcal{W}_\mu\left(\mathbb{R}_+;\mathcal{H},\cV_{R,M}^\rb\right)$, for some $\mu>0$, such that for all $t>0$ it holds that
\begin{equation}\label{eq:semi_discrete_rom}
\begin{aligned}
	\dotp{
		\partial^2_t
		u_{R,M}^{\normalfont\textrm{(rb)}}(t)
	}{
		v_{R,M}^{\normalfont\textrm{(rb)}}
	}_{\mathcal{H}}
	+ 
    \dual{\sfA u_{R,M}^\rb(t)}{v_{R,M}^\rb}_{\cV^\star\times\cV} = \dotp{f(t)}{v_{R,M}^\rb}_\cH
	% \mathsf{a}
	% \dotp{
	% 	u_{R}^{\normalfont\textrm{(rb)}}(t)
	% }{
	% 	v_{R}^{\normalfont\textrm{(rb)}}
	% }
	% = 
	% \dotp{f(t)}{v_{R}^{\normalfont\textrm{(rb)}} }_{\mathcal{H}}, 
	\quad
	\forall v_{R,M}^{\normalfont\textrm{(rb)}} \in \cV_{R,M}^\rb,
\end{aligned}
\end{equation}
equipped with vanishing initial conditions, i.e., 
$u_{R,M}^\rb(0)=0$ and $\partial_t u_{R,M}^\rb(0)=0$. 
\end{problem}

\begin{equation}
\begin{aligned}
	\dotp{
		\partial^2_t
		u_{R,M}^{\normalfont\textrm{(rb)}}(t)
	}{
		v_{R,M}^{\normalfont\textrm{(rb)}}
	}_{\mathcal{H}}
	+ 
    \dual{\sfA u_{R,M}^\rb(t)}{v_{R,M}^\rb}_{\cV^\star\times\cV} = \dotp{f(t)}{v_{R,M}^\rb}_\cH
	% \mathsf{a}
	% \dotp{
	% 	u_{R}^{\normalfont\textrm{(rb)}}(t)
	% }{
	% 	v_{R}^{\normalfont\textrm{(rb)}}
	% }
	% = 
	% \dotp{f(t)}{v_{R}^{\normalfont\textrm{(rb)}} }_{\mathcal{H}}, 
	\quad
	\forall v_{R,M}^{\normalfont\textrm{(rb)}} \in \cV_{R,M}^\rb,
\end{aligned}
\end{equation}

\begin{equation}\label{eq:min_laplace_domain}
	\mathcal{V}^{\text{(rb)}}_R
	=
	\argmin_{
	\substack{
		\mathcal{V}_R \subset \mathcal{V}
		\\
		\text{dim}(\mathcal{V}_R)\leq R	
	}
	}	
	\norm{
		\mathcal{L}\{\partial^2_t u\}
		-
		\mathsf{P}_{\mathcal{V}_R}
		\mathcal{L}\{\partial^2_t u\}
	}^2_{\mathscr{H}^2_\mu(\mathcal{V})},
\end{equation}

\begin{align}\label{eq:elliptic_projection}
	\dual{
		\mathsf{A}
		\Pi^{\text{(rb)}}_R v
	}{
		\vrb
	}_{\mathcal{V}^\star \times \mathcal{V}^\star}
	=
	\dual{
		\mathsf{A} v
	}{
		\vrb
	}_{\mathcal{V}^\star \times \mathcal{V}^\star},
	\quad\forall \vrb\in\cVrb.
\end{align}

\begin{align}\label{eq:cea_elliptic_projection}
    \norm{u - \Pi^{\text{(rb)}}_R u}_\cV 
    \leq 
    \sqrt{\frac{C_\mathsf{A}}{c_\mathsf{A}}}
    \inf_{\vrb\in\cVrb}\norm{u-\vrb}_\cV,
\end{align}

\begin{proposition}\label{eq:best_approximation}
Let $u \in \mathcal{W}_\mu(\mathbb{R}_+;\mathcal{H},\mathcal{V})$ for some $\mu>0$ be the solution to \eqref{eq:wave}--\eqref{eq:ic2} 
and let $u^{\normalfont\text{(rb)}}_R$ be the solution to Problem~\ref{pr:sdpr}
with $\cVrb$ as in \eqref{eq:min_laplace_domain}.

Then, it holds that
\begin{equation}
\begin{aligned}
	\norm{
		u
		-
		u^{\normalfont\text{(rb)}}_R
	}_{L^2(\mathfrak{I};\mathcal{V})}
    +
    &
	\norm{
		\partial_t 
		\left(
			u
			-
			u^{\normalfont\text{(rb)}}_R
		\right)
	}_{L^2(\mathfrak{I};\mathcal{H})}
	\\
	&
    \lesssim
    \sqrt{\frac{C_\mathsf{A}}{c_\mathsf{A}}}
   \max\left\{T,T^2, \frac{T}{\min\{1,\sqrt{ c_{\mathsf{A}} }\}}\right\}
    \norm{
        \left(
            \normalfont\text{Id}
            -
            \mathsf{P}^{\normalfont\text{(rb)}}_R
        \right) \partial^2_t u
    }_{L^2(\mathfrak{I};\mathcal{V})},
\end{aligned}
\end{equation}
where $\mathsf{P}^{\normalfont\text{(rb)}}_R: \mathcal{V} 
\rightarrow \mathcal{V}^{\normalfont\text{(rb)}}_R$ denotes
the $\mathcal{V}$-orthogonal projection onto $\mathcal{V}^{\normalfont\text{(rb)}}_R$.
% \MS{Check $\min$ in previous ineq.}
\end{proposition}

\begin{proof}
For a.e. $t \in \IR_{+}$ define 
\begin{equation}\label{eq:def_eta}
	\eta^{\text{(rb)}}_R(t) 
	\coloneqq 
	u^{\normalfont\text{(rb)}}_R(t)
	   -
	\left(\Pi^{\text{(rb)}}_R u\right)(t)
	\in
	\IV^{\text{(rb)}}_R.
\end{equation} 
By subtracting \eqref{eq:wave} and \eqref{eq:semi_discrete_rom},
and recalling that by definition $\cVrb\subset\cV$ we obtain
\begin{equation}\label{eq:subs_rb}
	\dual{
		\partial^2_t
		\left(
			u^{\normalfont\text{(rb)}}_R(t)
			-
			u(t)
		\right)
	}{
		v^{\text{(rb)}}_R
	}_{\mathcal{V}^\star\times \mathcal{V}}	
	+
	\mathsf{a}
	\dotp{
		u^{\normalfont\text{(rb)}}_R(t)
		-
		u(t)
	}{
		v^{\text{(rb)}}_R
	}
	=
	0,
    \quad
    \forall v^{\text{(rb)}}_R\in \IV^{\text{(rb)}}_R.
\end{equation}
Therefore, it follows from \eqref{eq:subs_rb} and the definition
of the elliptic projection, stated in \eqref{eq:elliptic_projection}, that
$\eta^{\text{(rb)}}_R(t)$ satisfies
for a.e.~$t \in \mathfrak{I}$
\begin{equation}\label{eq:residual_eta}
\begin{aligned}
	\dual{	
		\partial^2_t
		\eta^{\text{(rb)}}_R(t)
	}{
		v^{\text{(rb)}}_R
	}_{\mathcal{V}^\star\times \mathcal{V}}	
    +
	\mathsf{a}
	\dotp{
		\eta^{\text{(rb)}}_R(t)
	}{
		v^{\text{(rb)}}_R
	}
	=
   	&
	\dual{
		\left(
			\text{Id}
			-
			\Pi^{\text{(rb)}}_R
		\right)
		\partial^2_t
		 u(t) 
	}{
		v^{\text{(rb)}}_R
	}_{\mathcal{V}^\star\times \mathcal{V}}	,	
    \quad
    \forall v^{\text{(rb)}}_R \in \IV^{\text{(rb)}}_R,
\end{aligned}
\end{equation}
equipped with vanishing initial conditions.
Set $v^{\text{(rb)}}_R = \partial_t \eta^{\text{(rb)}}_R(t) \in  \IV^{\text{(rb)}}_R$, thus
\begin{equation}
\begin{aligned}
    \frac{1}{2}
    \partial_t
    \left(
    	\dotp{	
    		\partial_t
    		\eta^{\text{(rb)}}_R(t)
    	}{
    	   \partial_t 
           \eta^{\text{(rb)}}(t)
    	}_{\mathcal{H}}
        +
        \mathsf{a}
        \dotp{
            \eta^{\text{(rb)}}_R(t)
        }{
            \eta^{\text{(rb)}}_R(t)
        }
    \right)
    =
    \dotp{
		\left(
			\text{Id}
			-
			\Pi^{\text{(rb)}}_R
		\right)
		\partial^2_t
		 u(t) 
	}{
		\partial_t \eta^{\text{(rb)}}_R(t)
	}_{\mathcal{V}^\star\times \mathcal{V}}.
\end{aligned}
\end{equation}
Integrating over $[0,s]$, for any $s \in  \mathfrak{I}$ and $\epsilon>0$, we have
\begin{equation}
\begin{aligned}
   \frac{1}{2} \norm{	
        \partial_t
        \eta^{\text{(rb)}}_R(s)
    }^2_{\mathcal{H}}
	+
	c_{\mathsf{A}}
	\norm{
		\eta^{\text{(rb)}}_R(s)
	}^2_{\mathcal{V}}
	\text{d}\tau
    &
	\leq
    \norm{
    \left(
        \text{Id}
        -
        \Pi^{\text{(rb)}}_R
    \right) \partial^2_t u
    }_{L^2(\mathfrak{I};\mathcal{H})}
	\norm{
		\partial_t \eta^{\text{(rb)}}_R
	}_{L^2(\mathfrak{I};\mathcal{H})}
    \\
    &
    \leq
    \frac{\epsilon}{2}
    \norm{
    \left(
        \text{Id}
        -
       \Pi^{\text{(rb)}}_R
    \right) \partial^2_t u
    }^2_{L^2(\mathfrak{I};\mathcal{H})}
	+
    \frac{1}{2\epsilon}
    \norm{
		\partial_t \eta^{\text{(rb)}}_R
	}^2_{L^2(\mathfrak{I};\mathcal{H})}.
\end{aligned}
\end{equation}
Integrating over $\mathfrak{I}$, we obtain
\begin{equation}
   \frac{1}{2}  \norm{	
        \partial_t
        \eta^{\text{(rb)}}_R
    }^2_{L^2(\mathfrak{I};\mathcal{H})}
	+
	c_{\mathsf{A}}
	\norm{
		\eta^{\text{(rb)}}_R
	}^2_{L^2(\mathfrak{I};\mathcal{V})}
    \leq
    \frac{\epsilon T}{2}
    \norm{
    \left(
        \text{Id}
        -
       \Pi^{\text{(rb)}}_R
    \right) \partial^2_t u
    }^2_{L^2(\mathfrak{I};\mathcal{H})}
	+
    \frac{T}{2\epsilon}
    \norm{
		\partial_t \eta^{\text{(rb)}}_R
	}^2_{L^2(\mathfrak{I};\mathcal{H})}.
\end{equation}
Thus with $\epsilon=2T$, we obtain
\begin{equation}\label{eq:bound_eta_rb}
    \norm{	
        \partial_t
        \eta^{\text{(rb)}}_R
    }_{L^2(\mathfrak{I};\mathcal{H})}
	+
	\norm{
		\eta^{\text{(rb)}}_R
	}_{L^2(\mathfrak{I};\mathcal{V})}
    \leq
    \frac{\sqrt{2}T}{\min\{\frac{\sqrt{2}}{2} ,\sqrt{ c_{\mathsf{A}}}\}}
    \norm{
        \left(
            \text{Id}
            -
            \Pi^{\text{(rb)}}_R
        \right) \partial^2_t u
    }_{L^2(\mathfrak{I};\mathcal{H})}.
\end{equation}
Therefore, the triangle inequality implies that
\begin{equation}
\begin{aligned}
	\norm{
		u
		-
		u^{\normalfont\text{(rb)}}_R
	}_{L^2(\mathfrak{I};\mathcal{V})}
    +
    &
	\norm{
		\partial_t 
		\left(u
		-
		 u^{\normalfont\text{(rb)}}_R
		 \right)
	}_{L^2(\mathfrak{I};\mathcal{H})}
	\\
    \leq
    &
	\norm{
		\eta^{\text{(rb)}}_R
	}_{L^2(\mathfrak{I};\mathcal{V})}
    +
    \norm{	
        \partial_t
        \eta^{\text{(rb)}}_R
    }_{L^2(\mathfrak{I};\mathcal{H})}
    \\
    &
    +
    \norm{
        \left(
            \text{Id}
            -
             \Pi^{\text{(rb)}}_R
        \right)  u
    }_{L^2(\mathfrak{I};\mathcal{V})}
    +
    \norm{
        \left(
            \text{Id}
            -
             \Pi^{\text{(rb)}}_R
        \right)  \partial_t u
    }_{L^2(\mathfrak{I};\mathcal{H})}.
\end{aligned}
\end{equation}
Observe that since $u(0) = \partial_t u(0) = 0$, and using the vector-valued version of
Poincaré's inequality, we obtain
\begin{equation}
    \norm{
    \left(
        \text{Id}
        -
         \Pi^{\text{(rb)}}_R
    \right)
    u}_{L^2(\mathfrak{I};\mathcal{V})}
    \leq
    T
    \norm{
    \left(
        \text{Id}
        -
        \Pi^{\text{(rb)}}_R
    \right)
    \partial_t u}_{L^2(\mathfrak{I};\mathcal{V})}
\end{equation}
and
\begin{equation}
    \norm{
        \left(
            \text{Id}
            -
             \Pi^{\text{(rb)}}_R
        \right)
        \partial_t u
    }_{L^2(\mathfrak{I};\mathcal{V})}
    \leq
    T
    \norm{
    \left(
        \text{Id}
        -
         \Pi^{\text{(rb)}}_R
    \right)
    \partial^2_t u}_{L^2(\mathfrak{I};\mathcal{V})},
\end{equation}
which shows that
\begin{equation}
    \norm{
        \left(
            \text{Id}
            -
            \Pi^{\text{(rb)}}_R
        \right)
        u
    }_{L^2(\mathfrak{I};\mathcal{V})}
    \leq
    T^2
    \norm{
    \left(
        \text{Id}
        -
         \Pi^{\text{(rb)}}_R
    \right)
    \partial^2_t u}_{L^2(\mathfrak{I};\mathcal{V})}.
\end{equation}
Recalling \eqref{eq:bound_eta_rb} and the continuity of the embedding
$\mathcal{V} \subset \mathcal{H}$, we obtain
\begin{equation}
\begin{aligned}
	\norm{
		u
		-
		u^{\normalfont\text{(rb)}}_R
	}_{L^2(\mathfrak{I};\mathcal{V})}
    +
    &
	\norm{
		\partial_t 
		\left(u
		-
		 u^{\normalfont\text{(rb)}}_R
		 \right)
	}_{L^2(\mathfrak{I};\mathcal{H})}
	\\
	&
    \lesssim
    \max
    \left\{
        T,T^2,\frac{\sqrt{2}T}{\min\{\frac{\sqrt{2}}{2} ,\sqrt{ c_{\mathsf{A}}}\}}
    \right\}
    \norm{
        \left(
            \text{Id}
            -
        		\Pi^{\text{(rb)}}_R
        \right) \partial^2_t u
    }_{L^2(\mathfrak{I};\mathcal{V})}.
\end{aligned}
\end{equation}
Together with \eqref{eq:cea_elliptic_projection}, this yields the final result. 
\end{proof}

\end{document}
